\section{Proof details and reliable numerical endpoints}\label{app:tech}

\subsection*{Eliminating the positive current at fixed rows}

For the row relation in \eqref{eq:rows}, singular values of the restricted positive Green operator can be small. A padded inverse with an arbitrary ridge would change the optimization. For a given material tuple $x$ set $Z=D_x-T_x C$. If $Z$ lies in the range of the positive-row map $F_+$, its minimum positive Gram is exactly $Z^*(F_+F_+^*)^\dagger Z$. Otherwise the local tuple is impossible. For a necessary relaxation one may drop this range condition and take an outward upper envelope of $F_+F_+^*$, adding the appropriate range treatment or a positive floor. A trace, directional ellipsoid or eigenvalue bound on $C^*C$ may be intersected with these restrictions provided it holds for the full material hypothesis. A bound derived assuming $\|SJ-Y\|\le r$ can establish inconsistency of that ball. It cannot be inserted without qualification into a lower bound on $\delta(u;Y)$ outside the ball.

Two short checks delimit the rank argument. If $K$ is singular and $\Pi_{\ker K}L=0$, use $K^\dagger L$ in \eqref{eq:center}. The null coordinates contribute zero to completion. Subsequent local row equations must retain these free variables. If the range condition fails, the linear current term along the null directions cannot be discarded. Compression to the source directions annihilating $\Pi_{\ker K}L$ gives a lower bound on that part of the residual. For the main numerical comparisons an open spectral gap is checked on both sides instead.

A uniform receiver-exact fit for each separate coherent input need not hold jointly. For a collection of balances in spatial regions $b$, let $J_a=C_0a$ with $SC_0=Y$ and $w_j\in\ker S$. For any $P_j\succeq0$ use $\sum_j(a^*P_ja)w_jw_j^*$ as the current covariance for input $a$. The moment data fit exactly. With $D_b$ the indicator of a region, its sector and optical values receive the source-matrix increments
\[
 \sum_j[\|D_bw_j\|^2-u\Re(w_j^*D_bGw_j)]P_j,\qquad
 \sum_j\Im(w_j^*D_bGw_j)P_j .
\]
Satisfying the resulting matrix inequality and equality in each region constructs a feasible lifted current for every unit $a$ at once. Its energy is bounded by the largest eigenvalue of $C_0^*C_0+\sum_j\|w_j\|^2P_j$. In the two four-region examples of section~\ref{sec:experiments}, eight invisible directions suffice. The respective energy bounds are 8,500 and 7,000. A physical response in the same material box has energy above 8,886.62, so a scalar cap valid uniformly for the box and all unit inputs cannot eliminate those fits. Output-dependent caps on each complete error ball, $4.983\,10^6$ and $1.035\,10^6$, also admit them. These are statements about coherent relaxations. There is no assumption of strong duality for an arbitrary problem with several quadratic constraints \cite{beck}.

For the global sector and optical balances alone, a stronger constructive argument is possible. Write $q_r(j),q_o(j)$ for a unit input and its current balances. Suppose one receiver-null vector $w$ has $\rho=w^*Rw>0$ and $w^*(I-u\Herm G+t_0R)w<0$. Suppose $x=C_0a$ has $q_o(x)<0$ and $q_r(x)+t_0q_o(x)<0$ for all unit $a$. Choose $|\beta|^2=-q_o(x)/\rho$ and choose a phase of $\beta$ cancelling the optical cross term. Of the two opposite such phases, one has nonpositive sector cross term. Thus $x+\beta w$ satisfies both global balances and keeps the receiver fit. A covariance $\epsilon I+c_\epsilon ww^*$, with $c_\epsilon=(-q_o(x)-\epsilon\tr R)/\rho$, gives a relative Slater point for the lifted global problem for small $\epsilon>0$, since the sector remains strict. The coherent dual optimum is then zero \cite{boyd}. This is a data- and operator-dependent sufficient construction.

Conversely, a joint moment relaxation averages over currents. Each sample can spend a different negative source direction, so the rank of the average need not be bounded by $\nu$. In a five-cell reciprocal algebraic control with $R=I/20\succ0$, a four-point current distribution satisfies the averaged balances and has zero receiver variance, while a single negative direction leaves the positive squared-distance bound $2317/2400$. The physical squared distance is $800/401$. Conversely a positive source matrix multiplier can strengthen the scalar-current construction. In a two-cell zero-propagation control its lower bound is $88/47$, while every scalar-weight comparison in that example is at most $373/200$. These exact controls are given with their matrices in the numerical supplement. Positive-semidefinite lifting and pointwise interlacing exploit different restrictions.

\subsection*{A center term that changes an exclusion}

For a passive reciprocal algebraic example take $u=\mu=1$, $\lambda=0$ and
\[
 G=\begin{pmatrix}2+\mathrm i/2&1/4\\1/4&\mathrm i/2\end{pmatrix},\quad
 S^*S=\begin{pmatrix}1/2&1/4\\1/4&1/2\end{pmatrix},\quad
 B=\begin{pmatrix}-3/4-2\mathrm i\\5/4+9\mathrm i/2\end{pmatrix}.
\]
Choose $Y=S(1,-3+6\mathrm i)^\top$, generated by $X=2I$. The hypothesis is the entire real square $[0,1]^2$. Then $K=\diag(-1/2,3/2)$, $J_0=(5/4-\mathrm i,-5/12+7\mathrm i/2)^\top$, $(E_0)_1=103/48-5\mathrm i/2$, and $H=199/48$. There is one source and one negative direction. At cell one, $\omega_1=24$ gives $\omega_1|f_1(x)|^2=x^2$. For $c=1$ and $M_Y=3/2$ the first leading minor of \eqref{eq:comp} is $204x^2-192x+47\ge31/17$. Its determinant is $(4367x^2-3768x+864)/96>0$, since its discriminant in the numerator is negative. Therefore $\delta^2\ge127/48$ or $35.0802\%$ of the signal norm in distance. The rank bound is zero at these fixed weights. Uniformly bounding the norm of the center separately from the negative-row gain gives a subtraction exceeding 73 and again gives zero. This passive algebraic example does not use an open-domain Helmholtz geometry.

\subsection*{Whole-current and projected dissipation}

Split the negative term in \eqref{eq:started} as $\|W_c(J-J_0)\|^2+\|W_f(J-J_0)\|^2$, interpreted as source Grams. Let $W_f$ have $m$ rows. For $\Gamma=W_cM^{-1}W_c^*$ and $Z_0=W_c(J_c-J_0)$, completion at $\tau>\|\Gamma\|_2$ gives, for every $J$,
\begin{multline*}
 \tau(J-J_c)^*M(J-J_c)
 -(J-J_0)^*W_c^*W_c(J-J_0)\\
 \succeq-\tau Z_0^*(\tau I-\Gamma)^{-1}Z_0 .
\end{multline*}
Subtract $\tau C_B$ and add to the original quadratic. The fully shifted form dominates $H-\mathcal B_\tau-(J-J_0)^*W_f^*W_f(J-J_0)$ for every free $J$. Its free residual epigraph, with $\mathsf D\succeq0$, therefore has minimum trace at least $\Tail_m(H-\mathcal B_\tau)$. For nonsingular shifted state matrix, rank interlacing attains the free minimum by completion of squares. Thus feeding the same scalar into the whole dissipation constraint weakly dominates the partial matrix cap at that scalar. The rank cap can be cheaper to evaluate when projected quantities are already available. A singular transition is handled by a nondegenerate neighboring weight and a limit of lower bounds. Source-dependent multipliers in \eqref{eq:prices} escape this comparison.

For the strict comparison with all common dissipative prices, let $B_\tau=Z_Y^*(R_S+\tau I)^{-1}Z_Y-C_B$. It decreases with $\tau$. Hence for $0\le a\le\tau\le b$, $\tr[H(\tau)]_+\le b\,\tr[B_a]_+$. We interval-evaluated this expression on 101 intervals through an endpoint greater than $6.70676$ for $t=0.2$, $u=1$. At the actual right endpoint $B_\tau\preceq0$ and remains so. For $t=0.05$, $u=0.75$, the same calculation uses 92 intervals to $3.73262$, beyond which the positive trace vanishes. This proves the upper bound over a continuum of prices used in section~\ref{sec:experiments}. The lower bounds in coherent bases use intervals about the exact unitary polar factor of the stored basis. The receiver primal upper endpoints use explicit contractions. These computations bound an optimization in receiver space. A feasible receiver point can have no material preimage.

\subsection*{Native grid arithmetic, work and continuum scope}

The discrete convolution is a compression $G=V^*\diag(g)V$ of a $192^2$ padded Fourier multiplier by an isometry $V$. For scalar weights the Hermitian state form is the compression of a real symbol plus $S^*S/\mu$. Bounding the omitted symbol from below by a scalar floor leaves a positive diagonal and signed finite-rank Fourier updates. A small signed Gram congruence bounds the state negative count from above. An interval negative subspace of the full native operator gives the matching lower count. A failure of a floor to settle a gap is retained as unresolved. For quadrant-reflection invariant weights a separate check decomposes the reflection-averaged Hermitian operator into four parity blocks, each of size 2,304. In each block a low-rank update $WW^*$ raises the trial negative modes above zero. A Cholesky residual, including bounds for assembling the matrix, proves a strictly positive floor for the updated block. Interlacing bounds the number below that floor by $\rank W$. The norm error of reflection averaging, including the receiver Gram and stored kernel, is subtracted before combining the blocks. Independent negative Rayleigh subspaces on the original grid supply the reverse count. Symmetry is used only for the state weights and array. The material still varies independently at every pixel. For the all-cell test, $X=(u/2)I+\Delta$ with $\|\Delta\|\le u/2$ yields
\[
 \|(I-XG)z\|^2\le
 2\|(I-uG/2)z\|^2+\tfrac{u^2}{2}\|Gz\|^2 .
\]
The compressed Fourier norm is bounded by the corresponding uncompressed norm. The same calculation now applies to a symbol with the additional subtraction $\omega[2|1-ug/2|^2+u^2|g|^2/2]$. At the selected weight 441 retained Fourier modes and the receiver Gram suffice to put the full-box negative count at most thirteen. Trial residuals are evaluated by ball arithmetic on the original grid, including the cross term in \eqref{eq:trial}.

For any row block $A$ the positive Gram is $A K_+^\dagger A^*$. On the positive space take a block solution trial $Z_+$ and residual $\Delta=\Pi_+ A^*-K Z_+$. With $g_+>0$ a verified positive spectral gap,
\[
 A K_+^\dagger A^*\preceq A Z_+ + Z_+^* A^* - Z_+^* K Z_+
                              +g_+^{-1}\Delta^*\Delta .
\]
This follows by completing a square with exact remainder $\Delta^*K_+^\dagger\Delta$. Projection errors, repeated eigenvector clusters, the stationary residual and both centers must be enclosed. We reorthogonalized the low eigenspace and checked its leakage into the opposite space. The upper leakage at cap one was $2.2\,10^{-11}$. Endpoint stacks at $x=0,u$ give a uniform sufficient frame using convexity of the squared norm of an affine row. For the center-dependent block, candidate univariate summands were checked on each whole interval by coefficient comparison with $[I;xI]^*P[I;xI]+x(u-x)Q$ for positive matrices $P,Q$. Inexact polynomial coefficients, trial positivity and the center terms enter one-sided error allowances. Small semidefinite proposals and iterative eigenpairs are not themselves error enclosures. Final checks use real and complex ball arithmetic \cite{arb}.

The optical and stationary operators are shared by all observations with the same array, cap and weights. The main cap-one center and inertia run took 121 s wall time. Correcting the projected basis took another 21 s, and the full one-sided 104-row Gram and sparse bounds took 87 s. A separate proof for the dense supports took 12 s after 19 s of proposal time. The independent whole-domain source calculation took 159 s including verification. For the two edge weights, checking all four reflection subspaces and the six source matrices took 311 s wall time and 298 s CPU in total. Across all eighteen physical endpoint checks the inversion proposals used 426 native full-wave evaluations and 591,441 floating-point FFT columns. The final interval evaluations used 576 additional columns and 146 CPU seconds. These are measured costs for different outputs. Counting only a final source-space diagonalization would hide most of the work.

For clarity about the continuum statement \eqref{eq:weyl}, localize to a larger bounded region and replace the singular part of $H_G$ there by the positive model $k^2(1-\Delta_{\R^2})^{-1}$. Its difference has lower order than $-2$, since the outgoing resolvent has the same high-frequency principal symbol and the remainder away from the diagonal is smooth. The signed eigenvalue counting function of this localized remainder is $o(\varepsilon^{-1})$ \cite{ponge}. Compression cannot increase its singular values. For the leading operator $A\ge0$, compression to $\Omega$ has the same nonzero spectrum as $A^{1/2}{\bf1}_\Omega A^{1/2}$. When $|\partial\Omega|=0$, smooth squared cutoffs below and above the indicator sandwich this in the operator order, with integrals tending to $|\Omega|$. This includes a square. The compact-support negative-order Weyl law gives $k^2\int\varphi^2/(4\pi\varepsilon)$ for each such cutoff. The variational counting inequalities with a fractional threshold for the lower-order part give the same leading constant after that fraction tends to zero. Finally decompose $C$ into a finite-rank part and a remainder of norm below any $\epsilon>0$. Interlacing places the count between those of $H_G$ at thresholds $(1\pm\epsilon)/u$, up to a fixed finite number of modes. Taking $u\to\infty$ first proves \eqref{eq:weyl}. This derivation supplies no uniform discretization error or frequency-dependent remainder estimate.

If $\|Y-Y_{\rm model}\|_F$ includes additive uncertainty $a$, gain error $\eta_g<1$, and a separately justified loaded-scene discrepancy $b$, these enter the distance comparison in receiver units. For diagonal gains of operator norm error at most $\eta_g$, a conservative budget from a measured field norm is $(a+b+\eta_g\|Y\|_F)/(1-\eta_g)$ when $a,b$ are additive errors in the observed coordinates. An empirical empty-scene discrepancy supplies a calibration scale. A deterministic claim about a loaded target requires a bound appropriate to that target and the admitted physical models.

